% Generated by blueprint/scripts/atlas.py --remarks from atlas-mapping.json.
% Edit the reviewed mapping, then regenerate this appendix.
\chapter{The expert atlas and its Lean correspondence}
\label{chap:atlas-correspondence}

The expert atlas contributes a mathematical reading order and 102 named nodes.
Its Lean blocks are uncompiled sketches. The following remarks record the
reviewed correspondence, including interfaces and related specializations.
They do not add proof-completion marks. The interactive atlas computes usage
and solid edges from the compiled theorem; dashed edges are expert connections
without a supporting path between the selected versions in that root's closure.
The external inputs and their actual project dependencies are framed together.

\begin{remark}[Expert node W2CWA]
\label{rem:atlas-W2CWA}
\lean{Kakeya.Assouad.WZ2PaperBodyConvexWolffBound}
WZ2 Definition 2.5: Normalized convex-containment count.
The count is normalized by the indexed family cardinality. The lower tube-count consequence is valid; no upper cardinality bound follows. The Lean predicate uses extended nonnegative real volume and a finite error constant.
See \ref{def:cwa}, \ref{rem:cwa-consequences}, \ref{rem:wz1-extremality}.
\end{remark}

\begin{remark}[Expert node W2Unit]
\label{rem:atlas-W2Unit}
\lean{Kakeya.Assouad.WZ2PaperAssouadUnitRescalingData}
WZ2 Definition 2.8 / Sec. 6.1: Geometric and shaded unit rescaling.
The public definition uses the actual circumscribed John ellipsoid and full affine images. Re-cubulation and replacement by standard tubes are later constructions with separate containment and volume comparisons; the two operations are not identified by definition.
See \ref{def:unit-rescaling}, \ref{rem:john-normalization}, \ref{rem:common-grains}.
\end{remark}

\begin{remark}[Expert node W2MW]
\label{rem:atlas-W2MW}
\lean{Kakeya.Assouad.WZ2PaperPureCWAAtNearbyScales}
WZ2 Definition 2.12: Convex Wolff axioms at every scale.
This is the public nearby-scale predicate, including essential distinctness and complete containment fibres. Only the requested scale is bounded by one; the selected scale may exceed one.
See \ref{def:all-scale-cwa}, \ref{rem:full-fibres}, \ref{rem:nearby-scale}.
\end{remark}

\begin{remark}[Expert node W2L2\_13]
\label{rem:atlas-W2L2_13}
WZ2 Lemma 2.13: Coarse and fine multi-scale inheritance.
This is a mathematical inheritance interface. No standalone proof of the displayed general inheritance lemma is linked. The closed structural theorem reconstructs its particular covers and fibres; the expert inheritance arrows do not establish arbitrary refinement closure.
See \ref{lem:multiscale-cwa}, \ref{rem:refinement-loss}.
\end{remark}

\begin{remark}[Expert node W2Sigma]
\label{rem:atlas-W2Sigma}
\lean{Kakeya.Assouad.PureWZ2Admissible,Kakeya.Assouad.WZ2PaperPureIsExtremal,Kakeya.Assouad.pure_wz2_node02_critical_extraction}
WZ2 Section 6.2: The counterexample deficit sigma\_CW.
Lean defines admissibility globally and chooses its supremum locally inside critical extraction. There is no separately exported constant sigma\_CW identified with the directional WZ1 exponent.
See \ref{def:extremal}, \ref{lem:critical-exponent}, \ref{rem:wz1-extremality}.
\end{remark}

\begin{remark}[Expert node W2D6\_1]
\label{rem:atlas-W2D6_1}
\lean{Kakeya.Assouad.PureWZ2Admissible,Kakeya.Assouad.WZ2PaperPureIsExtremal}
WZ2 Definition 6.1: Definition of epsilon-CW-extremal.
The Lean extremal configuration stores nonemptiness, nearby-scale CWA, aggregate shading density, and the volume upper bound. Its critical exponent is supplied explicitly.
See \ref{def:extremal}, \ref{rem:wz1-extremality}.
\end{remark}

\begin{remark}[Expert node W2Exists]
\label{rem:atlas-W2Exists}
\lean{Kakeya.Assouad.pure_wz2_node02_critical_extraction}
WZ2 Section 6.2 / WZ1 Lem. 3.1: Extremizers exist at arbitrarily small scales.
Existence at arbitrarily small scales is a field of the critical package produced by the extraction theorem. The same theorem also supplies positivity, a subunit ceiling, and the lower envelope; this is one shared Lean declaration, not three independent proofs.
See \ref{lem:critical-exponent}, \ref{rem:wz1-extremality}.
\end{remark}

\begin{remark}[Expert node W2P6\_2]
\label{rem:atlas-W2P6_2}
\lean{Kakeya.Assouad.PureWZ2PropStickyStatement,Kakeya.Assouad.pure_wz2_node03_prop_sticky}
WZ2 Proposition 6.2: Balanced extremal coarse and fine families.
The closed version is in the localized cubical model and retains both cardinality factors. Spatial balance is equality of fine-union volumes in retained coarse cubes, not merely uniform child counts.
See \ref{prop:sticky-structure}, \ref{rem:balanced-cover}, \ref{rem:refinement-loss}.
\end{remark}

\begin{remark}[Expert node L3\_3]
\label{rem:atlas-L3_3}
\lean{Kakeya.Assouad.PureWZ2PropStickyStatement,Kakeya.Assouad.pure_wz2_node03_prop_sticky}
WZ1 Lemma 3.3 / CW: Local rescaling and multiplicity control.
The expert card is a CW transfer with the factor s squared times child cardinality. The original WZ1 fibre lemma in the manuscript is directional. The related WZ2 structure theorem supplies balanced rescaled fibres; it is not a separately proved copy of the expert local interface.
See \ref{lem:wz1-rescaled-fibre}, \ref{prop:sticky-structure}, \ref{rem:wz1-extremality}.
\end{remark}

\begin{remark}[Expert node L3\_4]
\label{rem:atlas-L3_4}
\lean{Kakeya.Assouad.WZ1RobustCloseCountFromBalancedConclusion,Kakeya.Assouad.wz1_robust_close_count_from_balanced}
WZ1 Lemma 3.4 / CW: Robust transversality.
The WZ1 conclusion is an interface in the directional uniform-structure model. A separate proved implication assumes a coarse close-direction packing input and balanced data. This does not independently prove the expert CW refinement statement.
See \ref{lem:robust-transversality}, \ref{rem:wz1-extremality}.
\end{remark}

\begin{remark}[Expert node L3\_5]
\label{rem:atlas-L3_5}
WZ1 Lemma 3.5 / CW: Mild anisotropic rescaling preserves extremality.
The expert card combines restriction, affine transport, re-cubulation and restoration of CWA. Those obligations are distributed across the structural proof. A generic mild-rescaling theorem of this displayed scope is not assigned a proof-completion mark.
See \ref{lem:mild-rescaling}, \ref{rem:common-grains}.
\end{remark}

\begin{remark}[Expert node W2Zero]
\label{rem:atlas-W2Zero}
\lean{Kakeya.Assouad.PureWZ2Theorem5_2Statement,Kakeya.Assouad.PureWZ2Theorem5_2Unconditional,Kakeya.Assouad.pure_wz2_node08_theorem5_2}
WZ2 Section 6.5: contradiction: The multi-scale convex Wolff deficit is zero.
The code closes a contradiction from a positive critical package and returns Theorem 5.2. It does not export the expert equality sigma\_CW = 0 as a separate scalar theorem.
See \ref{thm:sticky}, \ref{rem:endpoint}.
\end{remark}

\begin{remark}[Expert node W2T5\_2]
\label{rem:atlas-W2T5_2}
\lean{Kakeya.Assouad.PureWZ2Theorem5_2Statement,Kakeya.Assouad.PureWZ2Theorem5_2Unconditional}
WZ2 Theorem 5.2: Discrete union estimate for multi-scale convex Wolff tubes.
The public theorem has no external proof arguments. Nonemptiness is explicit, density is aggregate, and a support ball is not an additional hypothesis.
See \ref{thm:sticky}, \ref{rem:endpoint}, \ref{rem:aggregate-density}.
\end{remark}

\begin{remark}[Expert node Dplane]
\label{rem:atlas-Dplane}
\lean{Kakeya.Assouad.WZ1PlaneMapData}
WZ2 Section 6.1: plane map: Plane map.
The plane map is defined on the shaded union. Values at geometric cell centres require occupied-point representatives or an explicit extension. Affine transport acts on normals by the inverse transpose.
See \ref{def:plane-map}, \ref{rem:plane-domain}.
\end{remark}

\begin{remark}[Expert node T4\_2]
\label{rem:atlas-T4_2}
\lean{Kakeya.CV.MultilinearKakeyaThreeStatement,Kakeya.CV.multilinear_kakeya_three,Kakeya.Assouad.paper_cv_broad_set_estimate}
WZ1 Theorem 4.2: Endpoint multilinear Kakeya input.
The finite-segment multilinear theorem is proved but is not the version used by the sticky root. The root reaches the cropped full-line broad-set estimate, which supplies the mass bound needed for planiness.
See \ref{thm:multilinear-kakeya}.
\end{remark}

\begin{remark}[Expert node L4\_3]
\label{rem:atlas-L4_3}
\lean{Kakeya.Assouad.WZ1WeakPlaninessAssemblyConclusion}
WZ1 Lemma 4.3 / CW: Weak planiness.
The matching WZ1 declaration defines a proposition in the directional uniform-structure model. No proof of that full conclusion is supplied by the definition. The expert card transfers the argument to CW extremality; the closed WZ2 grain package uses its own constructions. Solid paths show only the actual compiled references.
See \ref{lem:weak-plane-map}.
\end{remark}

\begin{remark}[Expert node L4\_4]
\label{rem:atlas-L4_4}
\lean{Kakeya.Assouad.WZ1CoarsePlaneMapConclusion}
WZ1 Lemma 4.4 / CW: Coarsen weak planiness to a plane map.
The linked WZ1 interface records a coarse incidence tolerance of 5 times rho after direction alignment, rather than an exact unit tolerance. It also assumes the older directional uniform structure. The expert CW transfer is therefore not certified by the existence of this definition.
See \ref{lem:coarse-plane-map}.
\end{remark}

\begin{remark}[Expert node L4\_5]
\label{rem:atlas-L4_5}
\lean{Kakeya.Assouad.WZ1TwoScalePlaneMapConclusion}
WZ1 Lemma 4.5 / CW: Small oscillation on a fixed cube scale.
The matching WZ1 declaration defines a proposition in the directional uniform-structure model. No proof of that full conclusion is supplied by the definition. The expert card transfers the argument to CW extremality; the closed WZ2 grain package uses its own constructions. Solid paths show only the actual compiled references.
See \ref{lem:plane-oscillation}.
\end{remark}

\begin{remark}[Expert node C4\_6]
\label{rem:atlas-C4_6}
\lean{Kakeya.Assouad.WZ1ExtremalOneScalePlaneMapConclusion}
WZ1 Corollary 4.6 / CW: Small oscillation for nearby points.
The matching WZ1 declaration defines a proposition in the directional uniform-structure model. No proof of that full conclusion is supplied by the definition. The expert card transfers the argument to CW extremality; the closed WZ2 grain package uses its own constructions. Solid paths show only the actual compiled references.
See \ref{cor:plane-oscillation}.
\end{remark}

\begin{remark}[Expert node L4\_7]
\label{rem:atlas-L4_7}
\lean{Kakeya.Assouad.WZ1FinitePlaneMapIterationConclusion}
WZ1 Lemma 4.7 / CW: Iterate to a Lipschitz plane map.
The matching WZ1 declaration defines a proposition in the directional uniform-structure model. No proof of that full conclusion is supplied by the definition. The expert card transfers the argument to CW extremality; the closed WZ2 grain package uses its own constructions. Solid paths show only the actual compiled references.
See \ref{lem:finite-lipschitz}.
\end{remark}

\begin{remark}[Expert node L4\_8]
\label{rem:atlas-L4_8}
\lean{Kakeya.Assouad.WZ1LipschitzPlaneMapConclusion}
WZ1 Lemma 4.8 / CW: Existence of Lipschitz plane data.
The matching WZ1 declaration defines a proposition in the directional uniform-structure model. No proof of that full conclusion is supplied by the definition. The expert card transfers the argument to CW extremality; the closed WZ2 grain package uses its own constructions. Solid paths show only the actual compiled references.
See \ref{lem:lipschitz-plane-map}.
\end{remark}

\begin{remark}[Expert node DAD]
\label{rem:atlas-DAD}
\lean{Kakeya.Assouad.PureWZ2PaperADSet1,Kakeya.Assouad.IsADSet1}
WZ2 Section 6.1: ADset: Upper AD covering condition.
AD denotes a one-sided upper covering condition. The internal bounded-interval predicate and the paper-facing predicate have separate radius conventions; neither asserts a lower measure bound.
See \ref{def:ad-set}, \ref{rem:upper-ad}.
\end{remark}

\begin{remark}[Expert node L4\_9]
\label{rem:atlas-L4_9}
\lean{Kakeya.Assouad.WZ1LocalGrainCubeCountConclusion}
WZ1 Lemma 4.9 / CW: Grain count in tau-cubes.
The grain-count quotient is tau/rho to the power 1 minus sigma. The linked WZ1 conclusion is an interface in the directional model, not a proof of the expert CW transfer.
See \ref{lem:grain-cube}.
\end{remark}

\begin{remark}[Expert node L4\_10]
\label{rem:atlas-L4_10}
\lean{Kakeya.Assouad.WZ1LocalGrainIntervalConclusion}
WZ1 Lemma 4.10 / CW: Grain count in thin prisms.
The expert thin-prism step is presented in the manuscript through interval-localized normal projections. The linked WZ1 predicate records that output and is not itself its construction.
See \ref{lem:grain-interval}.
\end{remark}

\begin{remark}[Expert node L4\_11]
\label{rem:atlas-L4_11}
\lean{Kakeya.Assouad.WZ1LocalGrainOneScaleConclusion}
WZ1 Lemma 4.11 / CW: Fixed-mesh spacing at one grain scale.
The expert statement controls interval lengths at one fixed covering mesh. The existing WZ1LocalGrainOneScaleConclusion instead asks for a full IsADSet1 output. That stronger proposition is only an interface here; the fixed-mesh bound alone does not prove it. The manuscript now keeps these formulations distinct.
See \ref{lem:grain-ad-one-scale}, \ref{rem:fixed-mesh-ad}.
\end{remark}

\begin{remark}[Expert node L4\_12]
\label{rem:atlas-L4_12}
\lean{Kakeya.Assouad.WZ1EveryScaleLocalGrainConclusion,Kakeya.Assouad.pure_wz2_every_scale_local_grain}
WZ1 Lemma 4.12 / CW: AD spacing at every grain scale.
The expert route iterates the covering mesh. The repository also proves pure\_wz2\_every\_scale\_local\_grain, but that exported theorem is not reached from this root. Its existence must not turn the expert chain into solid Lean dependencies.
See \ref{lem:local-grains-all-scales}, \ref{rem:common-grains}.
\end{remark}

\begin{remark}[Expert node W2P6\_3]
\label{rem:atlas-W2P6_3}
\lean{Kakeya.Assouad.PureWZ2GrainsFromCriticalStatement,Kakeya.Assouad.pure_wz2_node04_grains}
WZ2 Proposition 6.3: Global and local grains with Lipschitz data.
Both grain conclusions belong to the same selected family and shading. The closed PureWZ2 theorem replaces the older WZ1 interfaces; dependency paths determine which internal constructions it actually uses.
See \ref{prop:lipschitz-grains}, \ref{rem:common-grains}, \ref{rem:plane-domain}.
\end{remark}

\begin{remark}[Expert node Dtrap]
\label{rem:atlas-Dtrap}
\lean{Kakeya.Assouad.WZ1VerticalTrapezoid}
WZ1 Definition 5.1: Vertical trapezoids and parents.
The geometric trapezoid type is implemented. Parent compatibility, occupied heights and supported corrections are additional data, not consequences of the type alone.
See \ref{def:trapezoid}.
\end{remark}

\begin{remark}[Expert node Iadgap]
\label{rem:atlas-Iadgap}
AD exponent gap: Why the expansion alternative is impossible.
This elementary exponent comparison is a mathematical explanation of the projection exclusion. No independent Lean theorem with this exact interface is asserted; its arithmetic can be inlined in the actual proof.
See \ref{lem:ad-expansion-gap}.
\end{remark}

\begin{remark}[Expert node T5\_2]
\label{rem:atlas-T5_2}
\lean{Kakeya.Assouad.WZ1ProjectionDichotomyConclusion}
WZ1 Theorem 5.2: Dot-product expansion or orthogonal lines.
The coefficient set concentrates near the perpendicular line through the origin; the endpoint set lies near an affine line. The linked general dichotomy is a proposition interface. The proved separated result has a weaker one-endpoint alternative.
See \ref{thm:wz1-projection}, \ref{rem:projection-orientation}, \ref{rem:projection-model}.
\end{remark}

\begin{remark}[Expert node L5\_3]
\label{rem:atlas-L5_3}
\lean{Kakeya.Assouad.WZ1LocallyLinearOneScaleData}
WZ1 Lemma 5.3: Local linearity from the two grain structures.
The linked WZ1 object stores a local affine-approximation witness; an inductive data declaration does not provide that witness. The direct-rich WZ2 construction retains additional local/global compatibility and its quantitative hypotheses.
See \ref{lem:locally-affine}.
\end{remark}

\begin{remark}[Expert node L5\_4]
\label{rem:atlas-L5_4}
\lean{Kakeya.Assouad.WZ1LocallyLinearOneScaleConclusion}
WZ1 Lemma 5.4 / CW: Piecewise linear approximation at one scale.
The exported WZ1 conclusion is an interface in the older extremal model. The WZ2 proof constructs its hierarchy through its own retained-height and common-shading data, rather than proving this whole interface first.
See \ref{lem:piecewise-affine}.
\end{remark}

\begin{remark}[Expert node L5\_5]
\label{rem:atlas-L5_5}
\lean{Kakeya.Assouad.WZ1LocallyLinearHierarchyData}
WZ1 Lemma 5.5 / CW: Compatible approximations at many scales.
A hierarchy data type is present, but a constructor with hierarchy fields as arguments is not an existence proof. Approximation, density and compatibility must be established on one refinement.
See \ref{lem:affine-hierarchy}.
\end{remark}

\begin{remark}[Expert node C5\_6]
\label{rem:atlas-C5_6}
\lean{Kakeya.Assouad.WZ1NestedSlopeCorrectionData}
WZ1 Corollary 5.6 / CW: A nested trapezoid tree for the slope graph.
The linked nested-correction data encode compatibility of parents and children. They are not a closed construction of the expert trapezoid hierarchy from extremality.
See \ref{cor:trapezoid-hierarchy}.
\end{remark}

\begin{remark}[Expert node L5\_8]
\label{rem:atlas-L5_8}
\lean{Kakeya.Assouad.WZ1SmoothSegmentExtensionStatement,Kakeya.Assouad.wz1_smooth_segment_extension}
WZ1 Lemma 5.8: Smooth extension of an affine segment.
The repository proves a smooth segment extension with one fixed cutoff constant. The support margin and both derivative costs are explicit.
See \ref{lem:segment-extension}, \ref{rem:smoothing}.
\end{remark}

\begin{remark}[Expert node P5\_7]
\label{rem:atlas-P5_7}
\lean{Kakeya.Assouad.WZ1MultiscaleSlopeCorrectionData}
WZ1 Proposition 5.7: A trapezoid tree produces a C2 approximation.
The linked WZ1 multiscale-correction object is data, not the theorem producing all of those corrections. The terminal derivative comparison involving the next scale is read only before the final index.
See \ref{prop:trapezoid-smoothing}, \ref{rem:smoothing}.
\end{remark}

\begin{remark}[Expert node W2P6\_4]
\label{rem:atlas-W2P6_4}
\lean{Kakeya.Assouad.PureWZ2C2GrainsFromCriticalStatement,Kakeya.Assouad.pure_wz2_node05_c2_grains}
WZ2 Proposition 6.4: Global grains with a C2 slope.
The root uses the direct-rich C2 assembly. The output family, shading, density, volume bound and both grain estimates share one combined-selection witness.
See \ref{prop:c2-grains}, \ref{rem:common-grains}.
\end{remark}

\begin{remark}[Expert node L6\_2]
\label{rem:atlas-L6_2}
WZ1 Lemma 6.2: A thin projection forces near orthogonality.
This paper lemma explains the projection-to-angle mechanism. No standalone declaration of its full expert formulation is linked; later slab arguments retain their own quantitative geometric data.
See \ref{lem:wz1-dense-segment}.
\end{remark}

\begin{remark}[Expert node W1SlabPrep]
\label{rem:atlas-W1SlabPrep}
\lean{Kakeya.Assouad.LargeSlopeTechnicalStatement}
WZ1 Lemma 6.3: Steps 1-3: Popular grains in a thin horizontal slab.
This is an unnumbered preparation inside the slope argument. The related large-slope statement does not independently construct the popular slice, occupied square or nonzero projected normal. These must be retained as actual witnesses.
See \ref{lem:large-slope}.
\end{remark}

\begin{remark}[Expert node W2SlabCount]
\label{rem:atlas-W2SlabCount}
\lean{Kakeya.Assouad.LargeSlopeTechnicalStatement}
WZ2 Lemma 6.6: new Step 4: Convex Wolff contradicts a popular thin slab.
The expert count concerns whole tubes contained in an enlarged convex slab. Shaded mass in a prism alone is insufficient. The related technical interface is not a proof of this grouped preparation-and-count step.
See \ref{lem:large-slope}.
\end{remark}

\begin{remark}[Expert node W2L6\_6]
\label{rem:atlas-W2L6_6}
\lean{Kakeya.Assouad.LargeSlopeTechnicalStatement}
WZ2 Lemma 6.6: The slope cannot be flat on a popular slab.
LargeSlopeTechnicalStatement is a proposition-valued interface, distinct from the closed WZ2 large-slope package. A declaration link to the former does not certify the generic technical lemma.
See \ref{lem:large-slope}.
\end{remark}

\begin{remark}[Expert node L6\_4]
\label{rem:atlas-L6_4}
WZ1 Lemma 6.4 / CW: A thick slab with comparable nonzero derivative.
The expert card transfers a WZ1 populated-height argument to CW extremality. The manuscript records the original directional formulation for comparison; the final normalized WZ2 output is supplied by Proposition 6.5.
See \ref{lem:wz1-comparable-derivative}.
\end{remark}

\begin{remark}[Expert node W2P6\_5]
\label{rem:atlas-W2P6_5}
\lean{Kakeya.Assouad.PureWZ2LargeSlopeFromCriticalStatement,Kakeya.Assouad.pure_wz2_node06_large_slope}
WZ2 Proposition 6.5, part 1 double-prime: Nonsingular C2 global grains.
The closed output includes the normalized derivative bounds together with local and global grains on the same shading. Transporting only the slope would not establish this package.
See \ref{prop:large-slope}, \ref{rem:common-grains}.
\end{remark}

\begin{remark}[Expert node W2L6\_7]
\label{rem:atlas-W2L6_7}
\lean{Kakeya.Assouad.PureWZ2SmallTwistedProjectionData,Kakeya.Assouad.pure_wz2_node07_small_twisted_projection}
WZ2 Lemma 6.7: small-image interface: Extremality produces a small twisted image.
The proved small-image exponent is sigma minus the loss. The printed plus sign is stronger and unnecessary. Parameter-Frostman preparation is supplied on the same retained configuration before the lower bound is applied.
See \ref{lem:small-twisted-projection}, \ref{rem:projection-sign}, \ref{rem:parameter-frostman}.
\end{remark}

\begin{remark}[Expert node Dtwist]
\label{rem:atlas-Dtwist}
\lean{Kakeya.Assouad.twistedProjection,Kakeya.Assouad.twistedCurve}
WZ1 Def. 6.1 / WZ2 Sec. 6.5: Twisted projection.
The map and four-parameter curve formula are defined separately. A supporting line and a particular unit-segment carrier are not interchangeable when converting geometric containment counts to parameter counts.
See \ref{def:twisted-projection}.
\end{remark}

\begin{remark}[Expert node W2TW]
\label{rem:atlas-W2TW}
\lean{Kakeya.Assouad.TubeWolffBound}
WZ2 Section 7: Tube Wolff: The weaker radius-squared tube count.
Ordinary Tube Wolff containment does not by itself count all carriers with nearby supporting-line parameters, because longitudinal positions may differ. The actual analytic theorem assumes indexed parameter nonconcentration separately.
See \ref{def:tube-frostman}, \ref{rem:parameter-frostman}.
\end{remark}

\begin{remark}[Expert node W2DKT]
\label{rem:atlas-W2DKT}
\lean{Kakeya.Assouad.DiscreteSet.IsKatzTao,Kakeya.Assouad.DiscreteSet.IsFrostman}
WZ2 Definition 7.2: Absolute Katz-Tao entropy bound.
Katz-Tao is an absolute count using (r/delta)\textasciicircum{}s. The reversed printed ratio is not used. It is distinct from normalized Frostman counting.
See \ref{def:kt-frostman}, \ref{rem:kt-quotient}.
\end{remark}

\begin{remark}[Expert node W2DFrost]
\label{rem:atlas-W2DFrost}
\lean{Kakeya.Assouad.DiscreteSet.IsKatzTao,Kakeya.Assouad.DiscreteSet.IsFrostman}
WZ2 Definition 7.3: Normalized Frostman entropy.
Frostman counting is normalized by the total size of the indexed set. Passing to a subset or a coarse family changes that normalization and requires a retention or uniform-fibre bound.
See \ref{def:kt-frostman}, \ref{rem:kt-quotient}.
\end{remark}

\begin{remark}[Expert node W2DUnif]
\label{rem:atlas-W2DUnif}
\lean{Kakeya.Assouad.ParameterSpacingUniformTreeData,Kakeya.Assouad.parameter_spacing_uniform_tree_of_base}
WZ2 Definition 7.8: Uniform occupied-cell counts.
The concrete ParameterSpacingUniformTreeData stores exact branching for the retained parameter points, together with its base, depth, atomization and retention bounds. Its construction is separately proved under quantitative base and scale hypotheses. This is a specialized parameter-tree implementation, not an identification of every generic uniform-set formulation.
See \ref{def:uniform-set}, \ref{rem:uniform-parameters}.
\end{remark}

\begin{remark}[Expert node Dcinematic]
\label{rem:atlas-Dcinematic}
\lean{Kakeya.Cinematic.IsCinematicFamily}
WZ1 Definition 7.1: Cinematic families of C2 functions.
The pairwise cinematic predicate controls differences, not absolute C2 norms. The family-uniform maximal estimate fixes an additional absolute C2 bound.
See \ref{def:cinematic-family}, \ref{rem:pyz-uniformity}.
\end{remark}

\begin{remark}[Expert node T7\_2]
\label{rem:atlas-T7_2}
\lean{Kakeya.Cinematic.WZ2UniformC2Input,Kakeya.Cinematic.wz2_uniformC2_input}
WZ1 Thm. 7.2 = WZ2 Thm. 7.13: Cinematic maximal estimate.
The proved input fixes a Katz-Tao constant, a neighbourhood dilation and an absolute C2 bound before the scale threshold. The full alpha/zeta PYZ statement is a separate unproved interface in this correspondence.
See \ref{thm:cinematic-l32}, \ref{rem:cinematic-normalization}, \ref{rem:pyz-uniformity}.
\end{remark}

\begin{remark}[Expert node L7\_3]
\label{rem:atlas-L7_3}
\lean{Kakeya.Assouad.CinematicFamilyFromSlopeStatement,Kakeya.Assouad.cinematic_family_from_slope}
WZ1 Lemma 7.3: The twisted-line family is cinematic.
The three-parameter cinematic-family theorem is proved, but this particular exported theorem is outside the sticky root's dependency closure. The root uses other parameter-family constructions. The theorem assumes f(0)=0 and constructs cinematic metric and doubling bounds; the absolute C2 bound is a further quantitative model estimate.
See \ref{lem:wz1-cinematic-family}, \ref{rem:wz1-slope-offset}.
\end{remark}

\begin{remark}[Expert node W2L7\_6]
\label{rem:atlas-W2L7_6}
\lean{Kakeya.Assouad.FrostmanToKatzTaoStatement,Kakeya.Assouad.frostman_to_katz_tao,Kakeya.Assouad.FrostmanToKatzTaoWeightedCardStatement,Kakeya.Assouad.frostman_to_katz_tao_weighted_card}
WZ2 Lemma 7.6: Frostman sets contain large Katz-Tao subsets.
The root uses the weighted-cardinality extraction, which does not assume separated representatives and retains the Frostman and logarithmic deficit. The separately proved separated version is unused by this root.
See \ref{lem:frostman-to-kt}, \ref{rem:replication-deficit}.
\end{remark}

\begin{remark}[Expert node W2L7\_9]
\label{rem:atlas-W2L7_9}
\lean{DiscretisedFurstenbergEstimate.BasicUniformization.basic_uniformization_dyadic}
WZ2 Lemma 7.9: Sparse-scale uniform refinement.
The nearby implemented uniformization is the OS dyadic theorem. Its number of levels and block lengths remain distinct. This correspondence does not equate the two papers' displayed constants.
See \ref{lem:uniform-set}, \ref{lem:os-uniformization}, \ref{rem:uniform-parameters}.
\end{remark}

\begin{remark}[Expert node W2L7\_10]
\label{rem:atlas-W2L7_10}
\lean{Kakeya.Assouad.ScaleProfileCrossingStatement,Kakeya.Assouad.scale_profile_crossing}
WZ2 Lemma 7.10: Zoom to Frostman fibres of a uniform set.
The linked scale-profile theorem proves the numerical crossing argument. The geometric selection, uniformity and local Frostman witnesses are separate inputs; the arithmetic lemma alone is not the whole spacing theorem.
See \ref{lem:spacing}, \ref{rem:os-losses}.
\end{remark}

\begin{remark}[Expert node W2L7\_11]
\label{rem:atlas-W2L7_11}
\lean{Kakeya.Assouad.ParameterBlockAmplificationStatement,Kakeya.Assouad.parameter_block_amplification}
WZ2 Lemma 7.11: Replicate a local Katz-Tao cluster.
The implemented three-parameter replication produces about q\textasciicircum{}\{-1\} copies and keeps the factor h\textasciicircum{}\{1-s\}, with h=delta/q. It does not assert saturated cardinality delta\textasciicircum{}\{-1\} after a lossy extraction.
See \ref{lem:block-replication}, \ref{rem:replication}, \ref{rem:replication-deficit}.
\end{remark}

\begin{remark}[Expert node W2ImageUniform]
\label{rem:atlas-W2ImageUniform}
\lean{Kakeya.Assouad.projected_fiber_os_preparation}
WZ2 Lemma 7.12: preparation: Uniformize the projected incidence set.
The implemented projected-fibre preparation retains incidence mass as well as area uniformity. A large target subset alone does not ensure a dense pulled-back shading.
See \ref{lem:projection-one-step}.
\end{remark}

\begin{remark}[Expert node W2Slice]
\label{rem:atlas-W2Slice}
\lean{Kakeya.Assouad.parameter_window_cluster_cardinality}
WZ2 Lemma 7.12: parameter slicing: A rich fixed-c slice of R4 parameters.
The formal implementation uses a finite parameter window and an explicit indexed four-parameter Frostman input. The expert inference from ordinary Tube Wolff alone is stronger than the actual interface.
See \ref{lem:freeze-line-parameter}, \ref{lem:projection-one-step}, \ref{rem:parameter-frostman}.
\end{remark}

\begin{remark}[Expert node W2JetCount]
\label{rem:atlas-W2JetCount}
\lean{Kakeya.Assouad.parameter_block_representative_cinematic_pullback}
WZ2 Lemma 7.12: cinematic assembly: Convert Katz-Tao triples to separated curves.
The implementation transports representative parameters and their shadings into the cinematic family. Separation, metric distortion and the replication deficit are hypotheses or proved fields, not automatic properties of a bare parameter set.
See \ref{lem:wz1-cinematic-family}, \ref{lem:projection-one-step}, \ref{rem:replication-deficit}, \ref{rem:cinematic-normalization}.
\end{remark}

\begin{remark}[Expert node W2AreaCompare]
\label{rem:atlas-W2AreaCompare}
\lean{Kakeya.Assouad.projected_fiber_os_cinematic_globalization}
WZ2 Lemma 7.12: Holder and translations: From cinematic norm to a local area ratio.
Holder acts on planar projected incidence mass, not three-dimensional shaded volume. The selected implementation globalizes through OS cells and cinematic comparison, with all retained mass and fixed constants recorded.
See \ref{lem:cinematic-area}, \ref{rem:projection-scales}.
\end{remark}

\begin{remark}[Expert node W2L7\_12]
\label{rem:atlas-W2L7_12}
\lean{Kakeya.Assouad.parameter_block_projection_uniformization_from_os}
WZ2 Lemma 7.12: One-step fine-to-coarse projected-area comparison.
The expert one-step inequality is a mathematical guide. The linked implementation is a parameter-block projection uniformization assembled from explicit arithmetic, mass-pruning, OS preparation and globalization inputs. These are not identical statements.
See \ref{lem:projection-one-step}, \ref{rem:projection-scales}.
\end{remark}

\begin{remark}[Expert node IcoarseTW]
\label{rem:atlas-IcoarseTW}
WZ2 proof operation: Uniform Tube-Wolff coarsening.
The normalized count is transported in parameter space using comparable nonempty fibres. Density, clipping at the height window and target-neighbourhood compatibility are separate geometric obligations. No single generic coarsening constructor is claimed.
See \ref{lem:uniform-parameter-coarsening}, \ref{rem:parameter-frostman}.
\end{remark}

\begin{remark}[Expert node W2Iterate]
\label{rem:atlas-W2Iterate}
\lean{Kakeya.Assouad.parameter_block_positive_window_assembly}
WZ2 Proposition 7.1: iteration: Telescope the ratios up to a coarse scale.
The expert finite telescoping argument is retained as mathematics. The root closes Section 7 through the parameter-block positive-window assembly; an absent reference to a separately named telescoping lemma must not be replaced by an invented solid edge.
See \ref{lem:projection-iteration}, \ref{rem:projection-iteration}.
\end{remark}

\begin{remark}[Expert node W2P7\_1]
\label{rem:atlas-W2P7_1}
\lean{Kakeya.Assouad.TwistedProjectionParameterFrostmanEstimateStatement,Kakeya.Assouad.pure_wz2_twisted_projection_parameter_frostman_lower_bound}
WZ2 Proposition 7.1: Twisted-image lower bound under Tube Wolff.
The proved specialization assumes a bounded vertical chart, analytic distinctness, an indexed four-parameter Frostman bound, an upper cardinality bound and f(0)=0. The broader Tube Wolff-only paper statement is not independently marked proved.
See \ref{thm:twisted-parameter}, \ref{prop:twisted-tube-wolff}, \ref{rem:parameter-frostman}.
\end{remark}

\begin{remark}[Expert node W2P6\_8]
\label{rem:atlas-W2P6_8}
\lean{Kakeya.Assouad.TwistedProjectionParameterFrostmanEstimateStatement,Kakeya.Assouad.pure_wz2_twisted_projection_parameter_frostman_lower_bound}
WZ2 Proposition 6.8: Twisted-image lower bound under convex Wolff.
CWA yields a geometric radius-squared tube count, but does not by itself establish all the parameter hypotheses of the proved lower bound for arbitrary unit-segment carriers. The generic expert CWA corollary is therefore a comparison node, not a second closed theorem.
See \ref{prop:twisted-cwa}, \ref{thm:twisted-parameter}, \ref{rem:parameter-frostman}.
\end{remark}

\begin{remark}[Expert node DFrost]
\label{rem:atlas-DFrost}
\lean{Kakeya.Assouad.DiscreteSet.IsFrostman}
WZ1 Definition 8.1: Normalized discretized Frostman sets.
The expert grid-volume convention and the implemented separated finite-point convention use normalized counts. The comparison requires discretization and its constants; they are not identical definitions.
See \ref{def:wz1-frostman}.
\end{remark}

\begin{remark}[Expert node Dhyp]
\label{rem:atlas-Dhyp}
\lean{Kakeya.Assouad.wz1HypergraphFiber,Kakeya.Assouad.WZ1UniformHypergraphDensity}
WZ1 Definitions 8.2-8.4: Hypergraph fibres and uniform density.
The implemented density includes every coordinate projection, including the empty and full coordinate sets. A graph restricted to active vertices still has to preserve these fibre inequalities.
See \ref{def:uniform-hypergraph}.
\end{remark}

\begin{remark}[Expert node Dthin]
\label{rem:atlas-Dthin}
\lean{HasMeasureThinTubes}
WZ1 Definition 8.5: Thin tubes for probability measures.
Thin tubes are a property of an ordered pair of measures and a retained relation. The reverse ordered pair is an additional hypothesis.
See \ref{def:thin-tubes}.
\end{remark}

\begin{remark}[Expert node DthinD]
\label{rem:atlas-DthinD}
\lean{Kakeya.Assouad.HasDiscreteThinTubes}
WZ1 Definition 8.5': Discretized thin tubes.
The discrete predicate counts retained incidences at radii above delta and normalizes by endpoint cardinality. A measure-level statement does not automatically hold at every discrete centre.
See \ref{def:discrete-thin-tubes}, \ref{rem:thin-tube-loss}.
\end{remark}

\begin{remark}[Expert node Dsep]
\label{rem:atlas-Dsep}
\lean{Kakeya.Assouad.WZ1StandardSeparation}
WZ1 Definition 8.6: Standard separation.
The implementation uses bounded finite charts, endpoint separation and separation of coefficients from zero. These chart hypotheses remain visible.
See \ref{def:standard-separation}.
\end{remark}

\begin{remark}[Expert node L8\_1]
\label{rem:atlas-L8_1}
\lean{Kakeya.Assouad.WZ1HypergraphRefinementStatement,Kakeya.Assouad.wz1_hypergraph_refinement}
WZ1 Lemma 8.1: Hypergraph refinement.
The finite hypergraph pruning theorem is proved, including all coordinate subsets and the retained-edge fraction. Its mathematical output is the density interface used subsequently.
See \ref{lem:hypergraph-refinement}.
\end{remark}

\begin{remark}[Expert node T8\_2]
\label{rem:atlas-T8_2}
\lean{Kakeya.Assouad.projection_energy_total}
WZ1 Theorem 8.2: A discretized Kaufman estimate with dense fibres.
The averaged projection-energy theorem is a supporting result. It is not, by itself, the full dense-fibre robust projection theorem on this expert card; the mass selections and energy-to-support conversion remain separate.
See \ref{thm:robust-kaufman}, \ref{thm:kaufman-energy}.
\end{remark}

\begin{remark}[Expert node L8\_3]
\label{rem:atlas-L8_3}
\lean{Kakeya.Assouad.WZ1SmoothedToDiscreteThinTubesStatement,Kakeya.Assouad.wz1_smoothed_to_discrete_thin_tubes}
WZ1 Lemma 8.3: Pass from measure thin tubes to grid thin tubes.
The smoothed-to-discrete theorem increases the exceptional fraction from c to 2c and uses an explicit enlarged counting constant. It does not preserve the expert/paper exceptional fraction unchanged.
See \ref{lem:unsmooth-thin-tubes}, \ref{rem:thin-tube-loss}.
\end{remark}

\begin{remark}[Expert node L8\_4]
\label{rem:atlas-L8_4}
\lean{Kakeya.Assouad.WZ1ThinTubesLargeDotProductStatement,Kakeya.Assouad.wz1_thin_tubes_large_dot_product}
WZ1 Lemma 8.4: Thin tubes force a large dot-product image.
The finite theorem has standard separation, normalized Frostman counts and a uniformly dense graph as hypotheses. Those conditions cannot be replaced by total edge density alone.
See \ref{lem:thin-tubes-dots}.
\end{remark}

\begin{remark}[Expert node P8\_6]
\label{rem:atlas-P8_6}
\lean{radial_bootstrapping_measure_thin_tubes}
WZ1 Proposition 8.6: Radial bootstrap for measures.
The closed measure bootstrap includes support in a normalized unit ball. This is the same Lean measure theorem associated with the external OSW input, rather than a new proof depending on a duplicate theorem.
See \ref{prop:radial-bootstrap}, \ref{rem:radial-support}.
\end{remark}

\begin{remark}[Expert node P8\_6p]
\label{rem:atlas-P8_6p}
\lean{radial_bootstrapping_measure_thin_tubes,Kakeya.Assouad.WZ1SmoothedToDiscreteThinTubesStatement,Kakeya.Assouad.wz1_discrete_to_smoothed_thin_tubes,Kakeya.Assouad.wz1_smoothed_to_discrete_thin_tubes,Kakeya.Assouad.wz1_finite_osw_one_step_iteration}
WZ1 Proposition 8.6': Radial bootstrap for grid sets.
The expert discrete bootstrap keeps the paper constants and exceptional fraction. The implemented conversion and finite iteration have explicit smoothing losses; the measure theorem does not itself prove this exact discrete statement.
See \ref{prop:radial-bootstrap}, \ref{lem:unsmooth-thin-tubes}, \ref{rem:thin-tube-loss}.
\end{remark}

\begin{remark}[Expert node C8\_7]
\label{rem:atlas-C8_7}
\lean{Kakeya.Assouad.wz1_finite_osw_one_step_iteration}
WZ1 Corollary 8.7: Iterate the bootstrap up to exponent 1.
The linked theorem iterates measure-valued thin tubes with an explicit bootstrap premise and the recurrence 3\textasciicircum{}N c. The expert card is discrete: smoothing and return to point counts impose additional losses, so that measure theorem does not certify the exact discrete constants.
See \ref{cor:thin-tube-iteration}, \ref{rem:thin-tube-loss}.
\end{remark}

\begin{remark}[Expert node L8\_8]
\label{rem:atlas-L8_8}
\lean{Kakeya.Assouad.WZ1QuarterThinTubesStatement,Kakeya.Assouad.wz1_quarter_thin_tubes}
WZ1 Lemma 8.8: Obtain a fixed 1/4-thin-tube seed.
The quarter-exponent seed is proved for separated finite Frostman sets. Both normalized ball counts and line nonconcentration hypotheses are retained.
See \ref{lem:quarter-thin-tubes}.
\end{remark}

\begin{remark}[Expert node P8\_5]
\label{rem:atlas-P8_5}
\lean{Kakeya.Assouad.wz1_line_nonconcentration_projection}
WZ1 Proposition 8.5: Strip nonconcentration forces expansion.
The proof uses the quarter seed, finite bootstrap and dot-product comparison with propagated constants. Its theorem type retains the parameter order and separated finite-set model.
See \ref{prop:line-nonconcentration}.
\end{remark}

\begin{remark}[Expert node L8\_10]
\label{rem:atlas-L8_10}
\lean{Kakeya.Assouad.two_ends_reduction}
WZ1 Lemma 8.10: Two-ends localization to a strip.
The implemented two-ends selection uses a near-maximal strip and a retained-mass margin. No attainment of an unrestricted supremum is assumed.
See \ref{lem:two-ends-strip}.
\end{remark}

\begin{remark}[Expert node L8\_11]
\label{rem:atlas-L8_11}
\lean{Kakeya.Assouad.wz1_anisotropic_frostman_rescaling}
WZ1 Lemma 8.11: Restore Frostman after anisotropic rescaling.
Anisotropic rescaling retains a substantial subset and comparable cell masses before renormalizing Frostman counts. Applying the affine map to the whole set is not the stated conclusion.
See \ref{lem:anisotropic-frostman}.
\end{remark}

\begin{remark}[Expert node L8\_12]
\label{rem:atlas-L8_12}
\lean{Kakeya.Assouad.wz1_frostman_coarsening}
WZ1 Lemma 8.12: Frostman coarsening after pigeonholing.
Coarsening uses a common dyadic mass class; its logarithmic retention loss changes the normalized Frostman constant.
See \ref{lem:frostman-coarsening}.
\end{remark}

\begin{remark}[Expert node L8\_13]
\label{rem:atlas-L8_13}
\lean{Kakeya.Assouad.WZ1StripLocalizationDichotomyStatement,Kakeya.Assouad.wz1_strip_localization_dichotomy_from_kaufman}
WZ1 Lemma 8.13: Comparable orthogonal strips or expansion.
The output concerns active vertices of the hypergraph, with the coefficient strip perpendicular through zero. The linked statement and its proved-from-inputs version retain the Kaufman/geometry inputs explicitly.
See \ref{lem:strip-localization}, \ref{rem:active-vertices}, \ref{rem:projection-orientation}.
\end{remark}

\begin{remark}[Expert node P8\_9]
\label{rem:atlas-P8_9}
\lean{Kakeya.Assouad.WZ1WellSeparatedProjectionUnionConclusion,Kakeya.Assouad.wz1_well_separated_projection_union}
WZ1 Proposition 8.9: Separated, uniform version of the dichotomy.
The closed result gives affine-line concentration in at least one endpoint class. The expert card asks for both independent endpoint classes, a stronger conclusion. The coefficient line passes through zero in the perpendicular direction.
See \ref{prop:wz1-separated}, \ref{rem:endpoint-union}, \ref{rem:projection-orientation}.
\end{remark}

\begin{remark}[Expert node T5\_2p]
\label{rem:atlas-T5_2p}
\lean{Kakeya.Assouad.WZ1ProjectionDichotomyConclusion,Kakeya.Assouad.WZ1WellSeparatedProjectionUnionConclusion,Kakeya.Assouad.wz1_well_separated_projection_union}
WZ1 Theorem 5.2': Frostman form, without separation or uniformity.
The expert general Frostman card has the reversed coefficient/endpoint orientation and simultaneous concentration in two endpoint sets. The implemented separated union alternative is weaker and has corrected orientation. No closed theorem of this full expert scope is asserted.
See \ref{thm:wz1-projection}, \ref{prop:wz1-separated}, \ref{rem:projection-orientation}, \ref{rem:endpoint-union}.
\end{remark}

\begin{remark}[Expert node XBCT]
\label{rem:atlas-XBCT}
\lean{Kakeya.CV.MultilinearKakeyaThreeStatement,Kakeya.CV.multilinear_kakeya_three,Kakeya.Assouad.paper_cv_broad_set_estimate,Kakeya.CV.polynomial_visibility_to_unitScale_multilinear_kakeya_of_factorization}
External [1] / [8]: Endpoint multilinear Kakeya.
The expert leaves multilinear Kakeya external. The repository proves the required dimension-three input using CV visibility and its topological and geometric ingredients. The arbitrary-dimensional paper theorem is not inferred from that specialization. The linked public polynomial application is conditional on its geometric inputs; the private unit-scale proof assembles those inputs before the broad-set estimate uses it.
See \ref{thm:multilinear-kakeya}, \ref{thm:cv-general}.
\end{remark}

\begin{remark}[Expert node XExt]
\label{rem:atlas-XExt}
\lean{Kakeya.Assouad.PureWZ2GrainsFromCriticalStatement,Kakeya.Assouad.pure_wz2_node04_grains,LipschitzOnWith.extend_real}
External [29] / Mathlib: Scalar Lipschitz extension.
The scalar Lipschitz extension theorem is a Mathlib boundary input. Its incoming references can be audited, while Mathlib's internal proof is outside the project traversal.
See \ref{prop:lipschitz-grains}.
\end{remark}

\begin{remark}[Expert node XOSW]
\label{rem:atlas-XOSW}
\lean{radial_bootstrapping_measure_thin_tubes}
External [35, Lemma 2.9]: Radial thin-tube bootstrap.
This external measure bootstrap is implemented in the repository. Its actual incidence, radial and additive-combinatorial dependencies are expanded inside the External frame; source numbering differs between the manuscripts.
See \ref{prop:radial-bootstrap}, \ref{rem:radial-support}.
\end{remark}

\begin{remark}[Expert node XPYZ]
\label{rem:atlas-XPYZ}
\lean{Kakeya.Cinematic.Theorem1_7UniformC2Statement}
External: Pramanik-Yang-Zahl: Cinematic maximal: absolute Katz-Tao input.
The full normalized PYZ statement is only a proposition interface in the current correspondence. The proved uniform-C2 L3/2 specialization and its actual tangency/counting dependencies are shown inside the External frame.
See \ref{thm:pyz}, \ref{rem:pyz-uniformity}.
\end{remark}

\begin{remark}[Expert node XFO]
\label{rem:atlas-XFO}
\lean{Kakeya.Assouad.FrostmanToKatzTaoStatement,Kakeya.Assouad.frostman_to_katz_tao,Kakeya.Assouad.FrostmanToKatzTaoWeightedCardStatement,Kakeya.Assouad.frostman_to_katz_tao_weighted_card}
External: Fässler-Orponen Lem. 3.13: Content-to-Katz-Tao extraction.
The expert invokes Hausdorff-content extraction as a black box. The root instead uses finite weighted-cardinality Frostman-to-Katz-Tao extraction. No independently proved content theorem is identified here.
See \ref{lem:content-to-kt}, \ref{lem:frostman-to-kt}.
\end{remark}

\begin{remark}[Expert node XOSuniform]
\label{rem:atlas-XOSuniform}
\lean{DiscretisedFurstenbergEstimate.BasicUniformization.basic_uniformization_dyadic}
External: Orponen-Shmerkin Lem. 2.15: Uniform refinement black box.
The implemented dyadic uniformization is recorded with its own theorem and scale conventions. The expert's paper numbering is provenance, not a Lean declaration name.
See \ref{lem:os-uniformization}, \ref{rem:uniform-parameters}.
\end{remark}

\begin{remark}[Expert node XSWdecomp]
\label{rem:atlas-XSWdecomp}
\lean{DirecretisedFurstenbergEstimate.FinalProof.multiscale_decomposition_main}
External: Shmerkin-Wang Lem. 5.21: Decompose the scale-entropy profile.
The linked OS multiscale decomposition is a related proved formulation, but that exported theorem is unused by this root. This does not certify a direct reference to the expert's Shmerkin-Wang black box.
See \ref{prop:os-multiscale}, \ref{rem:os-losses}.
\end{remark}

\begin{remark}[Expert node XSWmerge]
\label{rem:atlas-XSWmerge}
\lean{CombinatorialKaufman.MergingLemmas.merge_adjacent_epsilonSuperlinear}
External: Shmerkin-Wang Lem. 5.20: Merge adjacent entropy intervals.
The nearby implemented merging theorem requires equal chord slopes on two adjacent epsilon-superlinear intervals and returns their common slope. The expert black box merges decreasing dimensions by a weighted average, which is more general. These statements are related, not identified.
See \ref{lem:profile-decomposition}.
\end{remark}

\begin{remark}[Expert node Igeom]
\label{rem:atlas-Igeom}
\lean{Kakeya.DeltaTube,Kakeya.Assouad.WZ2PaperOrdinaryIsEssentiallyDistinct,Kakeya.Streamlined.Shading,Kakeya.Streamlined.Shading.IsLambdaDense,Kakeya.Streamlined.TubeFamily}
Geometry API: Lines, tubes, cubes, shadings.
The geometric API consists of indexed tube and shading data, preserving multiplicities. It is not a single theorem or a set of uncompiled Config declarations.
See \ref{def:tube}, \ref{def:shading}, \ref{rem:aggregate-density}.
\end{remark}

\begin{remark}[Expert node Ientropy]
\label{rem:atlas-Ientropy}
\lean{Kakeya.Assouad.PureWZ2PaperADSet1,Kakeya.Assouad.IsADSet1,Metric.externalCoveringNumber}
Entropy API: Finite covering numbers.
The underlying Mathlib external covering number takes values in the extended natural numbers and can be infinite. Boundedness, finiteness and coercion to real counting quantities are additional steps. The expert finite-entropy API groups these operations rather than supplying one theorem.
See \ref{def:ad-set}.
\end{remark}

\begin{remark}[Expert node Imass]
\label{rem:atlas-Imass}
\lean{Kakeya.Streamlined.Shading,Kakeya.Streamlined.Shading.IsLambdaDense,Kakeya.Assouad.PureWZ2BalancedCoverData,Kakeya.Streamlined.Shading.mass,Kakeya.Streamlined.Shading.pointMultiplicity}
Refinement API: Mass and constant multiplicity.
Mass, point multiplicity and spatial balance are different quantities. The grouped API does not allow total mass retention to replace per-fibre or per-tube counting hypotheses.
See \ref{def:shading}, \ref{def:balanced-cover}, \ref{rem:balanced-cover}, \ref{rem:refinement-loss}.
\end{remark}

\begin{remark}[Expert node Icover]
\label{rem:atlas-Icover}
\lean{Kakeya.Assouad.WZ2PaperPurePartitioningCover,Kakeya.Assouad.WZ2PaperPureFullFibersAreCUniform,Kakeya.Assouad.PureWZ2BalancedCoverData}
Cover API: Balanced coarse/fine configurations.
Full geometric containment fibres and doubled-fibre disjointness determine the parent map. Balanced shaded covers add a common positive fine-union volume in each retained coarse cell.
See \ref{def:partition-cover}, \ref{def:balanced-cover}, \ref{rem:full-fibres}, \ref{rem:balanced-cover}.
\end{remark}

\begin{remark}[Expert node Iaffine]
\label{rem:atlas-Iaffine}
\lean{Kakeya.Assouad.WZ2PaperAssouadUnitRescalingData}
Affine transport API: Jacobian, direction and grid transport.
Unit rescaling data and the later affine transport constructions have different roles. Plane normals use the inverse transpose, and induced shadings require determinant and re-cubulation comparisons.
See \ref{def:unit-rescaling}, \ref{lem:mild-rescaling}, \ref{rem:common-grains}.
\end{remark}

\begin{remark}[Expert node Icalc]
\label{rem:atlas-Icalc}
\lean{Kakeya.Assouad.WZ1SmoothSegmentExtensionStatement,Kakeya.Assouad.wz1_smooth_segment_extension}
Calculus API: C2 bounds and compactly supported patches.
The proved segment-extension theorem is one concrete calculus input. The expert API also includes finite supported-correction sums and their compatibility; the single lemma does not construct a whole smoothing hierarchy.
See \ref{lem:segment-extension}, \ref{rem:smoothing}.
\end{remark}

\begin{remark}[Expert node Iparams]
\label{rem:atlas-Iparams}
\lean{Kakeya.Assouad.pure_wz2_node02_critical_extraction}
Quantifier API: Small-parameter and finite-iteration contracts.
Small-parameter ordering and finite schedules are distributed through actual theorem binders and arithmetic helpers. No single theorem with all the expert API obligations is asserted.
See \ref{lem:projection-iteration}, \ref{lem:critical-exponent}, \ref{rem:projection-iteration}.
\end{remark}

\begin{remark}[Expert node Iwolffbalance]
\label{rem:atlas-Iwolffbalance}
WZ2 proof operation: Restriction and rebuilding of cover witnesses.
A retained index fraction gives the single-scale loss A/q. Every surviving containment fibre has its own retained fraction; arbitrary global deletion does not preserve all-scale CWA. The actual structural construction rebuilds its specific witnesses.
See \ref{lem:cwa-refinement}, \ref{lem:multiscale-cwa}, \ref{rem:refinement-loss}.
\end{remark}

\begin{remark}[Expert node Ienvelope]
\label{rem:atlas-Ienvelope}
\lean{Kakeya.Assouad.pure_wz2_node02_critical_extraction,Kakeya.Assouad.pure_wz2_non_admissible_spec}
WZ2 Section 6.2: consequence: Uniform lower envelope for admissible unions.
The non-admissibility lemma gives a uniform strict lower bound once an exponent is known not to be admissible. The critical extraction proves this premise above its locally chosen supremum; constant multiplicity then requires a separate mass comparison.
See \ref{lem:critical-exponent}, \ref{rem:wz1-extremality}.
\end{remark}
